\section{Literature Review}
\label{sec:literature}

\subsection{Retail and Demand Forecasting}
Retail forecasting has a long history as an applied field distinct from generic time-series analysis, owing to domain-specific complications such as promotions, price elasticity, new-product introduction, and stockouts~\cite{fildes2022retail,ord2017principles}. Large-scale empirical competitions have played a disproportionate role in establishing what actually works in practice. The M4 competition~\cite{makridakis2020m4} compared 61 forecasting methods across 100{,}000 series and found that combination methods and, later, hybrid statistical-machine-learning approaches outperformed both pure statistical and pure machine learning methods in isolation. The subsequent M5 competition~\cite{makridakis2022m5} specifically targeted retail sales data from a large retailer and confirmed that gradient-boosted tree methods were highly competitive on real retail series, while also highlighting the difficulty of forecasting intermittent, low-volume items accurately. These findings directly motivate the multi-model, automatically-adjudicated design adopted in this paper, rather than committing to a single model family a priori.

\subsection{Statistical Time-Series Forecasting}
The Box-Jenkins ARIMA methodology~\cite{boxjenkins1970} and its seasonal extension SARIMA remain foundational statistical approaches, modeling a series as a function of its own past values and past forecast errors after appropriate differencing. Exponential smoothing methods, including the Holt-Winters method~\cite{winters1960holtwinters}, model level, trend, and seasonal components multiplicatively or additively and remain highly competitive baselines, particularly for shorter series where more complex models are prone to overfitting~\cite{hyndman2018forecasting,makridakis1998forecasting}. More recently, additive regression frameworks such as Prophet~\cite{taylor2018prophet} were designed explicitly for business forecasting at scale, decomposing a series into trend, multiple seasonalities, and holiday effects, and are notable for degrading gracefully on noisy or irregular business data.

\subsection{Machine Learning for Time Series}
Framing time-series forecasting as a supervised regression problem over engineered lag, rolling-statistic, and calendar features allows conventional machine learning regressors to be applied directly. Random forests~\cite{breiman2001randomforest} and gradient-boosted trees, particularly the XGBoost implementation~\cite{chen2016xgboost}, have become standard strong baselines in this framing and were prominent among top-performing methods in the M5 competition~\cite{makridakis2022m5}. Instance-based methods such as $k$-nearest-neighbor regression~\cite{cover1967knn} and kernel methods such as support vector regression~\cite{cortes1995svm} are comparatively lightweight alternatives that can be competitive on smaller datasets where tree ensembles or neural networks have insufficient data to realize their capacity advantage.

\subsection{Deep Learning for Time Series}
Recurrent neural architectures, especially the long short-term memory (LSTM) network~\cite{hochreiter1997lstm,gers2000forget} and the gated recurrent unit (GRU)~\cite{cho2014gru}, were designed to model long-range temporal dependence while mitigating the vanishing-gradient problem of simple recurrent networks. Convolutional alternatives, including the temporal convolutional network~\cite{bai2018tcn}, achieve comparable or superior performance on some sequence-modeling benchmarks with better parallelization. Attention-based transformer architectures~\cite{vaswani2017attention} have subsequently been adapted specifically for long-sequence time-series forecasting, for example in the Informer architecture~\cite{zhou2021informer}, while specialized forecasting-native deep architectures such as DeepAR~\cite{salinas2020deepar} and N-BEATS~\cite{oreshkin2020nbeats} target probabilistic and interpretable forecasting respectively. A consistent finding across this literature, and a central empirical justification for the frequency-aware gating mechanism introduced in this paper, is that deep sequence models require substantially more training observations than statistical or tree-based methods to generalize reliably~\cite{lim2021survey}; when applied to short series they are prone to memorizing noise rather than learning genuine temporal structure.

\subsection{Hybrid and Meta-Learning Approaches}
A separate line of work combines statistical and machine-learning components directly, for example residual-correction hybrids in which a statistical model captures linear trend and seasonality while a neural network models the remaining nonlinear residual~\cite{zhang2003hybrid}. Meta-learning approaches for retail forecasting~\cite{ma2021metalearning} instead learn to select or weight among candidate models based on series characteristics. The system described in this paper adopts a related but simpler and more transparent strategy: every candidate model is always trained and evaluated identically via walk-forward cross-validation, and selection is performed by an explicit, auditable composite scoring rule rather than a learned meta-model, prioritizing transparency over marginal accuracy gains.

\subsection{Intermittent Demand Forecasting}
Standard forecasting methods perform poorly on sparse, intermittent demand series in which a large fraction of periods have zero recorded demand, a pattern common for slow-moving retail items and spare parts. Croston's method~\cite{croston1972} addresses this by separately smoothing the non-zero demand size and the inter-demand interval and forecasting their ratio as a constant rate. The Syntetos-Boylan approximation~\cite{syntetos2005sba} subsequently identified and corrected a systematic positive bias in Croston's original estimator, and is the variant adopted in this paper's implementation.

\subsection{Evaluation Methodology for Time Series}
Correctly evaluating time-series models requires respecting temporal ordering; naive random train-test splits leak future information into the training set and inflate reported accuracy. Walk-forward, expanding-window cross-validation is the accepted alternative~\cite{bergmeir2012cv}, and the number of folds used should reasonably reflect the amount of data available, since too few folds increases the risk that a model appears best due to chance performance on a single split~\cite{januschowski2020criteria}. Broader reviews of forecasting practice~\cite{petropoulos2022forecasting} additionally emphasize that a single point accuracy metric such as root-mean-squared error is an incomplete criterion for model selection in practice, since it does not capture computational cost, forecast stability, or the risk profile of a badly-calibrated model; this observation directly motivates the multi-criteria composite scoring approach described in Section~\ref{sec:methodology}.

\subsection{Explainable AI and Retrieval-Augmented Generation}
Post-hoc, model-agnostic explanation techniques such as SHAP~\cite{lundberg2017shap}, grounded in cooperative game-theoretic Shapley values, and LIME~\cite{ribeiro2016lime}, based on local surrogate modeling, are widely used to attribute a model's prediction to its input features. Tree-based models additionally expose a native, computationally inexpensive impurity- or gain-based feature-importance measure as a byproduct of training, which the system described in this paper surfaces directly for its gradient boosting and random forest models rather than computing a separate post-hoc attribution.

Independently, the emergence of large, instruction-following language models~\cite{brown2020gpt3,devlin2019bert} has introduced natural-language generation as a complementary explanation modality: rather than exposing numerical attributions alone, a system can narrate its reasoning in prose intelligible to a non-technical stakeholder. The central risk of this approach is hallucination, the generation of fluent but factually incorrect statements. Retrieval-augmented generation~\cite{lewis2020rag} mitigates this risk by conditioning language-model generation on retrieved, verifiable source documents rather than relying purely on the model's parametric knowledge. The explanation layer described in this paper applies this principle twice: qualitative business advice is grounded in a retrieved, curated retail-domain knowledge base, while all quantitative claims are grounded directly in the deterministic metrics computed by the forecasting pipeline itself, with explicit prompt-level instructions preventing the language model from re-deriving or contradicting labels the system has already computed.

\subsection{Positioning of This Work}
Relative to the reviewed literature, this paper does not propose a novel forecasting algorithm. Its contribution is architectural and systems-oriented: it integrates a broad, heterogeneous ensemble spanning all four model families discussed above behind a single automated pipeline, introduces an explicit and auditable robust multi-criteria selection rule in place of raw error-metric ranking, applies an empirically justified frequency-based gate to prevent unreliable deep-learning training on insufficient data, and couples the resulting forecasts to a retrieval-augmented natural-language explanation layer that is explicitly constrained to remain consistent with the system's own deterministic outputs.
